\documentclass[11pt]{article}
\usepackage[a4paper,margin=26mm]{geometry}
\usepackage{amsmath,amssymb,booktabs}
\newcommand{\F}{\mathbb F}
\title{The parameter restrictions in Table 1}
\author{A source-table calculation for Leemans--Vandenschrick}\date{}
\begin{document}\maketitle
The first case of Section~5.3 of the source has odd characteristic and
parameters $a,a'\notin\{0,1,-1\}$. Its four involutions preserve a
nonsingular conic. When both maximal rank-three subgroups are $S_4$ or
$A_5$, Table~1 records the possible characteristics for the three
adjacent orders. We recover every numerical cell below. This is a
necessary parameter calculation; it neither proves the intersection
property nor replaces the separately acquired classification of the
resulting groups.

Write $X$ for the characteristic-polynomial variable. Direct
multiplication of the four displayed matrices on physical page~16 gives
\[
 \begin{gathered}
 (X-1)(X^2-2aX+1),\quad (X-1)(X^2+CX+1),\quad (X-1)(X^2-2a'X+1),\\
 C=-aa'+a+a'+1.
 \end{gathered}
\]
Thus the source's $b$ in the last characteristic polynomial on physical
page~17 must be $a'$. For end order three, $a=-1/2$; for end order five,
\[
 a=\frac{-1+t}{4}\quad\hbox{or}\quad\frac{-1-t}{4},
 \qquad t^2=5.
\]
The first assertion follows from the two primitive cube roots; the second
from $2a=\zeta+\zeta^{-1}$ for a primitive fifth root. In the present
nondegenerate parameter case characteristic five is excluded from an
order-five end, since it would give $a=1$.

The adjacent order in the middle is three, four or five precisely when,
respectively,
\[
 C=1,\qquad C=0,\qquad C^2-C-1=0.                 \tag{1}
\]
For the semisimple cases this follows from the reciprocal eigenvalue
pair. In characteristic three or five the repeated eigenvalue one
requires a qualification: the two central involutions are distinct, so
their product is not the identity. Its nontrivial unipotent part in
three dimensions has order three or five, respectively. Their scalar
centres cannot shorten the order, because the matrices have eigenvalue
one. Thus (1) remains valid in the relevant characteristics.

An order-four middle cannot be adjacent to an order-five end in $S_4$
or $A_5$: the former contains no element of order five and the latter
none of order four. These are the dashes in the source table; they are
not omitted equations. An order-four middle with two order-three ends
would require $C=-1/4=0$, which is impossible.

For the other cells, substitute the three possible values of each end
parameter. Symmetry in $a,a'$ and replacement $t\mapsto-t$ reduce the
nine source rows to the following four forms. The last two columns are
possible characteristics; the actual equation in (1) must still hold
for the chosen root $t$ in the ambient field.
\[
\begin{array}{c|c|c|c}
 (a,a')&C&\text{middle }3&\text{middle }5\\\hline
 (-1/2,-1/2)&-1/4&5&11\\
 (-1/2,(-1\pm t)/4)&(1\pm3t)/8&\varnothing&59\\
 ((-1\pm t)/4,(-1\pm t)/4)&(1\pm5t)/8&19&3,11\\
 ((-1+t)/4,(-1-t)/4)&3/4&\varnothing&19
\end{array}
\]
For clarity, these restrictions follow by integer elimination rather
than a search over fields. Reduce the left sides of (1) modulo $t^2-5$
and clear powers of two. Their integer numerators $u_3,u_5$, in the same four rows and
using the plus choice of sign, are
\[
\begin{array}{c|r|r|r|r}
 &u_3&\operatorname{Res}_t(t^2-5,u_3)
 &u_5&\operatorname{Res}_t(t^2-5,u_5)\\\hline
 1&-5&25&-11&121\\
 2&3t-7&4&-9t-13&-236\\
 3&5t-7&-76&27-15t&-396\\
 4&-1&1&-19&361
\end{array}                                                     \tag{2}
\]
A common root in a field of odd characteristic forces that prime to
divide the corresponding integer resultant. Factoring (2) gives exactly
the displayed characteristic lists. Conversely, in each listed
characteristic the two polynomials have a common root in a finite
extension, and it does not force either end parameter into
$\{0,1,-1\}$. The supplied exact polynomial calculation checks those
small gcd assertions and retains all nine signed rows. The ambient
field must contain this common root; one may not assert that both
choices of $\sqrt5$ work merely from the prime.

This recovers the parameter restrictions in the source table and corrects
the last displayed characteristic polynomial.
Its title must not be read as a certificate that every listed parameter
choice gives a rank-four C-group. The actual generation and intersection
conditions are supplied by the separately reviewed conic classification.
The calculation adds no whole-original, conjecture, or measurement credit.
\end{document}
